\chapter{Trigonometría: el círculo unitario}\label{ch:g11:trigo}

La trigonometría comienza en triángulos rectángulos, pero su verdadero hogar es el
\emph{círculo unitario}, dónde \omterm{def:g11:trigo:cossin}{coseño} y \omterm{def:g11:trigo:cossin}{señor} convertirse funciones de un arbitrario
número real. Este capítulo instala el \omterm{def:g11:trigo:radian}{radián}, el devanado de lo real
línea alrededor del círculo, los valores exactos que debe saber de memoria y el
simetrías que generan todas las identidades clásicas. el estudio de
$\cos $ y $\sin $ como funciones --- variaciones, derivadas --- es llevado
afuera en \cref{ch:g12:trigo}.

\section{Radianes y el devanado}

\begin{definition}[Radián]\label{def:g11:trigo:radian}
Sea $\mathcal C $ ser el círculo de radio $1$ centrado en el origen $ O $,
y $ I(1, 0)$. El \emph{radián}\index{radián} medida de un ángulo en el
centro de $\mathcal C $ es la longitud del arco que corta $\mathcal C $.
Como el círculo completo tiene circunferencia $2\pi $:
\[
360^\circ = 2\pi \text{ rad},
\qquad
180^\circ = \pi \text{ rad},
\qquad
1 \text{ rad} = \frac{180^\circ}{\pi} \approx 57.3^\circ .
\]
\end{definition}

\begin{method}[Mudado]\label{met:g11:trigo:convert}
Grados y radianes son proporcionales: multiplicar por $\frac{\pi}{180}$ ir
de grados a radianes, por $\frac{180}{\pi}$ al revés. Por ejemplo
$60^\circ = 60 \times \frac{\pi}{180} = \frac{\pi}{3}$, y
$\frac{3\pi}{4} = \frac{3 \times 180^\circ}{4} = 135^\circ $.
\end{method}

\begin{definition}[Enrollando la línea en el círculo.]\label{def:g11:trigo:winding}
a cada uno número real $ t $, asocia el punto $ M(t)$ de $\mathcal C $
obtenido al caminar una distancia $\abs t $ a lo largo del círculo de $ I $,
en sentido antihorario si $ t \geq 0$, en el sentido de las agujas del reloj si $ t < 0$. desde el
la circunferencia es $2\pi $, los números $ t $ y $ t + 2k\pi $ ($ k \in \Z $)
llegar al mismo punto: $ M(t + 2k\pi) = M(t)$.
\end{definition}

\begin{example}
$ M(0) = I $; $ M\!\left(\frac{\pi}{2}\right) = (0, 1)$, la parte superior del
círculo; $ M(\pi) = (-1, 0)$; $ M\!\left(-\frac{\pi}{2}\right) = (0, -1)$;
y $ M\!\left(\frac{9\pi}{4}\right) = M\!\left(\frac{\pi}{4} +
2\pi\right) = M\!\left(\frac{\pi}{4}\right)$.
\end{example}

\section{Coseno y seno}

\begin{definition}[Coseno y seno]\label{def:g11:trigo:cossin}
por cada real $ t $, el \emph{coseño}\index{coseño} y
\emph{señor}\index{señor} de $ t $ son los \omterm{def:g10:coordgeom:system}{coordenadas} del punto $ M(t)$ de
el círculo unitario:
\[
M(t) = (\cos t,\ \sin t).
\]
\end{definition}

\begin{omfigure}
\begin{tikzpicture}[scale=2.0]
  \draw[->] (-1.3,0) -- (1.4,0) node[right] {$ x $};
  \draw[->] (0,-1.3) -- (0,1.4) node[above] {$ y $};
  \draw[gray] (0,0) circle (1);
  \coordinate (O) at (0,0);
  \coordinate (M) at (55:1);
  \draw[thick,omDef] (O) -- (M);
  \draw[dashed] (M) -- (M |- O) node[below, font=\small] {$\cos t $};
  \draw[dashed] (M) -- (M -| O) node[left, font=\small] {$\sin t $};
  \draw[->,omThm,thick] (0.35,0) arc (0:55:0.35);
  \node[omThm, font=\small] at (27:0.55) {$ t $};
  \fill (M) circle (0.7pt) node[above right] {$ M(t)$};
  \fill (1,0) circle (0.7pt) node[below right] {$ I $};
\end{tikzpicture}

{\small El círculo unitario: el número real $ t $, enrollado en sentido antihorario desde
$ I $, aterriza en el punto $ M(t)$ cuyo \omterm{def:g10:coordgeom:system}{coordenadas} \emph{definir} $\cos t $
y $\sin t $.}
\end{omfigure}

\begin{proposition}[Primeras propiedades]\label{prop:g11:trigo:first}
A pesar de $ t \in \R $ y $ k \in \Z $:
\[
-1 \leq \cos t \leq 1, \qquad
-1 \leq \sin t \leq 1, \qquad
\cos^2 t + \sin^2 t = 1,
\]
\[
\cos(t + 2k\pi) = \cos t, \qquad \sin(t + 2k\pi) = \sin t .
\]
\end{proposition}

\begin{proof}
$ M(t)$ se encuentra en el círculo unitario, por lo que su \omterm{def:g10:coordgeom:system}{coordenadas} están entre $-1$ y
$1$, y satisfacen la del círculo \omterm{def:g10:algebra:equation}{ecuación} $ x^2 + y^2 = 1$, que es el
identidad $\cos^2 t + \sin^2 t = 1$. Replanteos de periodicidad
$ M(t + 2k\pi) = M(t)$ (\cref{def:g11:trigo:winding}).
\end{proof}

\begin{example}\label{ex:g11:trigo:identity}
Si $\sin t = \frac35$ y $ t \in \intcc{\frac\pi2}{\pi}$ (segundo
cuadrante), entonces $\cos^2 t = 1 - \frac{9}{25} = \frac{16}{25}$, entonces
$\cos t = \pm\frac45$; en el segundo cuadrante el \omterm{def:g10:coordgeom:system}{abscisa} es negativo,
por lo tanto $\cos t = -\frac45$.
\end{example}

Los valores a saber de memoria:
\begin{center}
\begin{tabular}{c|ccccc}
$ t $ & $0$ & $\dfrac{\pi}{6}$ & $\dfrac{\pi}{4}$ & $\dfrac{\pi}{3}$ & $\dfrac{\pi}{2}$\\[6pt]
\hline\rule{0pt}{14pt}%
$\cos t $ & $1$ & $\dfrac{\sqrt3}{2}$ & $\dfrac{\sqrt2}{2}$ & $\dfrac12$ & $0$\\[6pt]
$\sin t $ & $0$ & $\dfrac12$ & $\dfrac{\sqrt2}{2}$ & $\dfrac{\sqrt3}{2}$ & $1$
\end{tabular}
\end{center}

\begin{remark}
Una ayuda para la memoria: el senos leer
$\frac{\sqrt0}{2}, \frac{\sqrt1}{2}, \frac{\sqrt2}{2}, \frac{\sqrt3}{2},
\frac{\sqrt4}{2}$, y el \omterm{def:g11:trigo:cossin}{coseños} son la misma lista invertida.
\end{remark}

\section{Ángulos asociados}

\begin{proposition}[Ángulos asociados]\label{prop:g11:trigo:associated}
\index{ángulos asociados}
A pesar de $ t \in \R $:
\[
\begin{aligned}
\cos(-t) &= \cos t, & \sin(-t) &= -\sin t,\\
\cos(\pi - t) &= -\cos t, & \sin(\pi - t) &= \sin t,\\
\cos(\pi + t) &= -\cos t, & \sin(\pi + t) &= -\sin t,\\
\cos\left(\tfrac{\pi}{2} - t\right) &= \sin t, &
\sin\left(\tfrac{\pi}{2} - t\right) &= \cos t .
\end{aligned}
\]
\end{proposition}

\begin{proof}
Cada línea expresa una simetría del círculo.

$ M(-t)$ es el reflejo de $ M(t)$ en el $ x $-eje: el \omterm{def:g10:coordgeom:system}{abscisa} es
sin cambios, el ordenada cambia de signo.

$ M(\pi - t)$ es el reflejo de $ M(t)$ en el $ y $-eje: caminar
hacia atrás desde $\pi $ por $ t $ aterriza opuesto (horizontalmente) a caminar
adelante desde $0$ por $ t $. El \omterm{def:g10:coordgeom:system}{abscisa} cambia de signo, el ordenada es
sin cambios.

$ M(\pi + t)$ es diametralmente opuesto $ M(t)$ (\omterm{def:g10:stats:mean}{media} vuelta): ambos
\omterm{def:g10:coordgeom:system}{coordenadas} cambiar de signo.

$ M\!\left(\frac\pi2 - t\right)$ es el reflejo de $ M(t)$ en el
linea diagonal $ y = x $, que intercambia los dos \omterm{def:g10:coordgeom:system}{coordenadas}.
\end{proof}

\begin{omfigure}
\begin{tikzpicture}[scale=2.0]
  \draw[->] (-1.35,0) -- (1.45,0) node[right] {$ x $};
  \draw[->] (0,-1.35) -- (0,1.45) node[above] {$ y $};
  \draw[gray] (0,0) circle (1);
  \fill (40:1) circle (0.7pt) node[above right, font=\small] {$ M(t)$};
  \fill (-40:1) circle (0.7pt) node[below right, font=\small] {$ M(-t)$};
  \fill (140:1) circle (0.7pt) node[above left, font=\small]
    {$ M(\pi - t)$};
  \fill (220:1) circle (0.7pt) node[below left, font=\small]
    {$ M(\pi + t)$};
  \draw[black, dashed, thin] (40:1) -- (-40:1);
  \draw[black, dashed, thin] (40:1) -- (140:1);
  \draw[black, dashed, thin] (140:1) -- (220:1);
  \draw[omDef, thin] (0,0) -- (40:1);
  \draw[omDef, thin] (0,0) -- (140:1);
  \draw[omDef, thin] (0,0) -- (220:1);
  \draw[omDef, thin] (0,0) -- (-40:1);
\end{tikzpicture}

{\small Los cuatro puntos asociados: $ M(-t)$ espejos $ M(t)$ en el
$ x $-eje, $ M(\pi - t)$ en el $ y $-eje, y $ M(\pi + t)$ es diametralmente
opuesto. Cada identidad de \cref{prop:g11:trigo:associated} se lee
esta foto.}
\end{omfigure}

\begin{example}
$\cos\frac{2\pi}{3} = \cos\left(\pi - \frac\pi3\right)
= -\cos\frac\pi3 = -\frac12$, y
$\sin\frac{7\pi}{6} = \sin\left(\pi + \frac\pi6\right)
= -\sin\frac\pi6 = -\frac12$. La tabla de cinco valores, ampliada por el
simetrías, cubre todo el círculo.
\end{example}

\section{Resolver ecuaciones trigonométricas}

\begin{theorem}[las ecuaciones $\cos t = \cos a $ y $\sin t = \sin a $]
\label{thm:g11:trigo:equations}
Para números reales $ t $ y $ a $:
\[
\cos t = \cos a \iff t = a + 2k\pi \ \text{or}\ t = -a + 2k\pi
\quad (k \in \Z),
\]
\[
\sin t = \sin a \iff t = a + 2k\pi \ \text{or}\ t = \pi - a + 2k\pi
\quad (k \in \Z).
\]
\end{theorem}

\begin{proof}
Dos puntos del círculo unitario tienen el mismo \omterm{def:g10:coordgeom:system}{abscisa} exactamente cuando son
iguales o reflejos entre sí en el $ x $-eje; por
\cref{prop:g11:trigo:associated}, el reflejo de $ M(a)$ es $ M(-a)$. Entonces
$\cos t = \cos a $ medio $ M(t) = M(a)$ o $ M(t) = M(-a)$, es decir.\
$ t = \pm a $ hasta un múltiplo de $2\pi $. De manera similar, igual ordenados \omterm{def:g10:stats:mean}{media}
puntos iguales o reflexiones en el $ y $-eje, y la reflexión de $ M(a)$
es $ M(\pi - a)$.
\end{proof}

\begin{method}[resolviendo $\cos t = c $ en un intervalo]\label{met:g11:trigo:solve}
\begin{enumerate}
  \item Encontrar \emph{uno} ángulo $ a $ con $\cos a = c $, de la tabla de
        valores conocidos.
  \item Escribe las soluciones generales. $ t = \pm a + 2k\pi $
        (\cref{thm:g11:trigo:equations}).
  \item Elige el números enteros $ k $ ese terreno dentro del solicitado \omterm{def:g10:numbers:interval}{intervalo},
        y enumere las soluciones.
\end{enumerate}
Proceder de la misma manera para $\sin t = c $ con $ t = a + 2k\pi $ o
$ t = \pi - a + 2k\pi $.
\end{method}

\begin{example}\label{ex:g11:trigo:solve}
Resolver $\cos t = \frac12$ en $\intoc{-\pi}{\pi}$. Una solución de
$\cos a = \frac12$ es $ a = \frac{\pi}{3}$. Las soluciones generales son
$ t = \frac\pi3 + 2k\pi $ y $ t = -\frac\pi3 + 2k\pi $. Adentro
$\intoc{-\pi}{\pi}$, solo $ k = 0$ contribuye:
$ t \in \left\{-\frac\pi3,\ \frac\pi3\right\}$.

Resolver $\sin t = -\frac{\sqrt2}{2}$ en $\intco{0}{2\pi}$. Un ángulo es
$ a = -\frac\pi4$; las soluciones son $ t = -\frac\pi4 + 2k\pi $ y
$ t = \pi + \frac\pi4 + 2k\pi = \frac{5\pi}{4} + 2k\pi $. En
$\intco{0}{2\pi}$: $ t = \frac{7\pi}{4}$ (de $ k = 1$) y
$ t = \frac{5\pi}{4}$.
\end{example}

\begin{omfigure}
\begin{tikzpicture}[scale=2.0]
  \draw[->] (-1.35,0) -- (1.45,0) node[right] {$ x $};
  \draw[->] (0,-1.2) -- (0,1.3) node[above] {$ y $};
  \draw[gray] (0,0) circle (1);
  \draw[omProp, thick] (0.5,-1.05) -- (0.5,1.15)
    node[above, font=\small] {$ x = \tfrac12$};
  \fill (60:1) circle (0.7pt)
    node[above right, font=\small] {$ M\!\left(\tfrac\pi3\right)$};
  \fill (-60:1) circle (0.7pt)
    node[below right, font=\small] {$ M\!\left(-\tfrac\pi3\right)$};
  \draw[omDef, thin] (0,0) -- (60:1);
  \draw[omDef, thin] (0,0) -- (-60:1);
\end{tikzpicture}

{\small resolviendo $\cos t = \frac12$: la línea vertical $ x = \frac12$
(naranja) corta el círculo unitario en dos puntos, simétrico en el $ x $-eje
--- de ahí las dos familias de soluciones $ t = \pm\frac\pi3 + 2k\pi $.}
\end{omfigure}

\section{Ejercicios}

\begin{exercise}[$\star $]\label{exo:g11:trigo:1}
Convertir a radianes: $30^\circ $, $45^\circ $, $120^\circ $, $270^\circ $.
Convertir a grados: $\frac{\pi}{5}$, $\frac{5\pi}{6}$, $\frac{7\pi}{4}$,
$\frac{2\pi}{9}$.
\end{exercise}

\begin{exercise}[$\star $]\label{exo:g11:trigo:2}
Coloque en el círculo unitario los puntos. $ M(t)$ para
\[
t = \frac{2\pi}{3},\quad
t = -\frac{\pi}{4},\quad
t = \frac{17\pi}{6},\quad
t = -\frac{7\pi}{2},
\]
después de reducir cada uno a un valor en $\intoc{-\pi}{\pi}$ módulo $2\pi $.
\end{exercise}

\begin{exercise}[$\star $]\label{exo:g11:trigo:3}
Usando la tabla y el \omterm{prop:g11:trigo:associated}{ángulos asociados}, proporcione los valores exactos de
\[
\cos\frac{3\pi}{4}, \quad
\sin\frac{5\pi}{6}, \quad
\cos\left(-\frac{\pi}{3}\right), \quad
\sin\frac{4\pi}{3}, \quad
\cos\frac{11\pi}{6} .
\]
\end{exercise}

\begin{exercise}[$\star $]\label{exo:g11:trigo:4}
Dado $\cos t = \frac{5}{13}$ con $ t \in \intoo{-\frac\pi2}{0}$, calcular
$\sin t $ exactamente.
\end{exercise}

\begin{exercise}[$\star\star $]\label{exo:g11:trigo:5}
Simplifica, para cualquier real $ x $:
\[
A = \cos(\pi + x) + \cos(-x) + \sin\left(\frac\pi2 - x\right) - \cos x,
\qquad
B = \sin(\pi - x) + \sin(\pi + x) + \sin(-x) .
\]
\end{exercise}

\begin{exercise}[$\star\star $]\label{exo:g11:trigo:6}
Resolver $\intoc{-\pi}{\pi}$:
\[
\cos t = \frac{\sqrt3}{2}, \qquad
\sin t = \frac12, \qquad
\cos t = -1 .
\]
\end{exercise}

\begin{exercise}[$\star\star $]\label{exo:g11:trigo:7}
Resolver $\intco{0}{2\pi}$:
\[
\sin t = -\frac{\sqrt3}{2}, \qquad
\cos t = 0, \qquad
2\sin t + 1 = 0 .
\]
\end{exercise}

\begin{exercise}[$\star\star $]\label{exo:g11:trigo:8}
Resolver $\R $ el \omterm{def:g10:algebra:equation}{ecuación} $\cos 2t = \cos t $. (Usar
\cref{thm:g11:trigo:equations} con $ a = t $y hablar sobre las dos familias
de soluciones.)
\end{exercise}

\begin{exercise}[$\star\star $]\label{exo:g11:trigo:9}
Muéstralo de verdad $ x $,
\[
\bigl(\cos x + \sin x\bigr)^2 + \bigl(\cos x - \sin x\bigr)^2 = 2 .
\]
\end{exercise}

\begin{exercise}[$\star\star $]\label{exo:g11:trigo:10}
una rueda de radio $30$ Rollos de cm sin deslizarse. ¿A través de qué ángulo (en
radianes, luego en grados) gira cuando la bicicleta avanza $1.5$ ¿metro?
¿Qué distancia avanza la bicicleta durante un giro completo de la rueda?
\end{exercise}

\begin{exercise}[$\star\star\star $]\label{exo:g11:trigo:11}
Resolver $\intoc{-\pi}{\pi}$ la desigualdad
$\cos t \leq \frac12$. (Dibuje el círculo unitario, marque la región donde
\omterm{def:g10:coordgeom:system}{abscisa} es como \omterm{def:g10:functions:extrema}{máximo} $\frac12$, y lea el \omterm{def:g10:numbers:interval}{intervalo} de valores de
$ t $.)
\end{exercise}

\section{Problema: La noria y la marea}

\begin{problem}[{Problema del fin de semana --- radianes como distancia recorrida,
y el círculo unitario como reloj maestro de cada periódico.
fenómeno}]
\label{pb:g11:trigo:1}
Una noria te lleva alrededor de un círculo; la marea lleva el
el agua del puerto arriba y abajo en el mismo círculo matemático. cada
fenómeno periódico --- ruedas, mareas, sonido, estaciones --- lee
su horario fuera del círculo unitario de este capítulo
(\cref{def:g11:trigo:cossin}), y el ecuaciones de
\cref{met:g11:trigo:solve} decir las horas exactas. este problema
convierte, enrolla, modela y resuelve --- y se establece en el camino por qué
Los matemáticos miden ángulos en radianes.

\medskip
\noindent\textbf{Part I --- radianes, the rolled distance.}
\begin{enumerate}
  \item Convertir a radianes: $30^\circ $, $135^\circ $,
        $300^\circ $; y en grados: $\frac{3\pi}{4}$,
        $\frac{5\pi}{6}$ (\cref{met:g11:trigo:convert}).
  \item Todo el punto de radianes: un arco de ángulo $ t $ (en
        radianes) en un círculo de radio $ r $ tiene longitud exactamente
        $ r\,t $. ¿Qué arco forma un ángulo de $2$ radianes cortar en un
        circulo de radio $3$ ¿metro? (No $\pi $ en cualquier lugar --- eso es
        el punto.)
  \item una rueda de radio $30$ rollos de cm $1$ km sin
        resbalando. ¿A través de cuántos \emph{radianes} ¿Se ha vuelto?
        y cuantas vueltas completas son?
  \item Colocar en el círculo unitario y evaluar exactamente:
        $\cos\frac{2\pi}{3}$, $\sin\left(-\frac{\pi}{4}\right)$,
        y $\cos\frac{19\pi}{6}$ (reducir módulo $2\pi $ primero;
        \cref{prop:g11:trigo:associated}).
  \item un angulo $ t \in \intoo{\frac\pi2}{\pi}$ satisface
        $\sin t = \frac35$. Usando
        $\cos^2 t + \sin^2 t = 1$ y el cuadrante, encuentre
        $\cos t $ y $\tan t $.
\end{enumerate}

\medskip
\noindent\textbf{Part II --- The Ferris wheel.}
Una rueda gigante tiene radio. $60$ m, su centro $65$ m por encima del
suelo y gira una vez cada $30$ minutos. Usted aborda en el
punto más bajo en el momento $ t = 0$ (en minutos), por lo que tu altura obedece
\[
h(t) = 65 - 60 \cos\!\left(\frac{\pi t}{15}\right).
\]
\begin{enumerate}[resume]
  \item Justifique la fórmula: verifique el ángulo girado después $ t $
        minutos, la altura a $ t = 0$, y explica el menos
        firmar.
  \item Calcula tu altura en $ t = 7.5$, $ t = 15$ y
        $ t = 20$ minutos (valores exactos).
  \item ¿A qué hora estás? \emph{primero} en $95$ ¿metro? (Resolver
        $ h(t) = 95$ con \cref{met:g11:trigo:solve}.)
  \item durante cual tiempo \omterm{def:g10:numbers:interval}{intervalo} del viaje ¿estás al menos?
        $95$ m de alto --- ¿y por cuánto tiempo en total? (Comparar
        \cref{exo:g11:trigo:11}.)
  \item ¿Dónde cambia la altura? \emph{lo más rápido}? comparar el
        subir durante el primer minuto
        ($ h(1) - h(0)$) con la subida entre minutos $7$ y
        $8$ y formular la observación (la herramienta que hace
        es exacto --- la derivada de la sinusoide --- llega
        en el grado 12).
  \item El turno del ingeniero: diseñar la fórmula de altura de una rueda
        con altura máxima $40$ m, altura mínima $4$ m y
        periodo $12$ minutos, embarque por la parte inferior a las
        $ t = 0$.
\end{enumerate}

\medskip
\noindent\textbf{Part III --- The tide and the circle's
symmetries.}
En un puerto, la profundidad del agua en metros, $ t $ horas después
medianoche, es
\[
d(t) = 7 + 3 \sin\!\left(\frac{\pi t}{6}\right).
\]
\begin{enumerate}[resume]
  \item Indique las profundidades máxima y mínima, los momentos en que
        aparecen por primera vez y el período de la marea.
  \item Un buque de carga necesita $8.5$ m de agua. Resolver
        $ d(t) \geq 8.5$ en $\intco{0}{12}$: durante qué ventana
        ¿Puede entrar y cuánto mide la ventana? ¿Cuándo
        siguiente ventana abierta?
  \item Vuelva a derivar dos identidades de \omterm{prop:g11:trigo:associated}{ángulos asociados} directamente de
        simetrías circulares --- $\sin(\pi - t) = \sin t $ (espejo
        a lo largo del eje vertical) y $\cos(\pi + t) = -\cos t $
        (\omterm{def:g10:stats:mean}{media} vuelta) --- y enuncie la identidad complementaria
        $\cos\left(\frac\pi2 - t\right) = \sin t $.
  \item paridad (\cref{pb:g11:func:1} vocabulario): clasificar
        $\sin $, $\cos $, y $ t \mapsto \sin t + t^3$ como par o
        \omterm{def:g11:func:parity}{impar}, con pruebas de una línea del círculo.
  \item Resolver $\intco{0}{2\pi}$:
        $2\sin^2 t - \sin t - 1 = 0$. (Una cuadrática en
        disfraz: factorizarlo como en \cref{pb:g11:quad:1}, entonces
        lee el círculo.)
\end{enumerate}

\medskip
\noindent\textbf{Part IV --- The mathematician's angle.}
\begin{enumerate}[resume]
  \item cuantos grados es $1$ \omterm{def:g11:trigo:radian}{radián}? Luego una trampa de calculadora:
        que es más grande, $\sin(1)$ (\omterm{def:g11:trigo:radian}{radián} modo) o
        $\sin(1^\circ)$ --- ¿y aproximadamente por qué factor? moraleja
        para la configuración de la calculadora?
  \item Resolver $\cos t = \sin t $ en $\intco{0}{2\pi}$, usando el
        círculo (¿dónde están? \omterm{def:g10:coordgeom:system}{abscisa} y ordenada ¿igual?).
  \item Para ángulos pequeños, el arco y su \omterm{def:g11:trigo:cossin}{señor} casi coinciden:
        comparar $\sin(0.1)$ con $0.1$ (cinco decimales), y
        dar el error relativo. Este "ángulo pequeño"
        \omterm{def:g10:numbers:approx}{aproximación}'' gobierna barcos y péndulos; el teorema
        detrás de él ($\frac{\sin t}{t} \to 1$) es una estrella de la
        Capítulo de límites de grado 12.
  \item Final --- el círculo unitario como reloj maestro, una frase
        por artículo: radianes convertir ángulos en números honestos (arco
        $=$ radio $\times $ ángulo); \omterm{def:g11:trigo:cossin}{coseño} y \omterm{def:g11:trigo:cossin}{señor} son los
        sombra y altura de la manecilla del reloj; las simetrias del circulo
        son las identidades; resolviendo $\cos t = c$ lee tiempos
        fuera del horario; y todo fenómeno periódico --- rueda,
        marea, sonido --- este reloj lleva un disfraz.
\end{enumerate}
\end{problem}
